\originalchapter{熵、条件信息与无噪声编码}
\sourcelecture{33}

连续投掷 $k$ 枚公平硬币有 $2^k$ 个等可能结果, 用 $k$ 个二进制位就能把结果写清楚. 如果结果不等可能, 经常出现的结果应使用短码, 少见结果可以使用长码. 本章把这一直觉变成两个可以证明的结论: 平均码长不能低于熵, 而适当编码可以使它与熵相差不到一个比特. 条件熵则说明已有信息怎样改变剩余的不确定性.

\originalsection{信息量与熵}
本章 $\log_2$ 始终表示以 $2$ 为底的对数, $\ln$ 表示自然对数. 概率为 $p>0$ 的结果的信息量定义为 $-\log_2p$. 两个独立结果的联合概率为乘积, 对数把乘积变为和, 因而信息量相加. 概率 $1/2$ 的结果携带一个比特, 概率 $1/8$ 的结果携带三个比特. 这不是说看到一个罕见结果就一定实际写了三个二进制位, 而是提供了平均编码长度的标尺.

\begin{definition}[离散熵]
设 $X$ 的取值集合有限或可数, 质量函数为 $p(x)$. 定义
\[
 H(X)=-\sum_xp(x)\log_2p(x)=\E[-\log_2p(X)].
\]
约定 $0\log_20=0$, 因为 $\lim_{t\downarrow0}t\log_2t=0$. 可数情形允许 $H(X)=+\infty$; 正概率结果的信息量均非负, 所以此级数没有正负抵消的歧义.
\end{definition}
熵依赖各结果的概率, 不依赖给结果起什么名称. 对概率为零的结果, $-\log_20$ 本身无穷, 但该结果不会发生, 在期望和中其贡献按上述连续延拓取零.

\begin{example}
计算常量、$m$ 点均匀分布、伯努利分布和正整数上的几何分布的熵.
\end{example}
\begin{solution}
常量的熵为 $-1\log_21=0$. $m$ 点均匀分布的熵为
$-m(1/m)\log_2(1/m)=\log_2m$. 若 $X\sim\Bin(1,p)$, 则
\[
 H(X)=h(p):=-p\log_2p-(1-p)\log_2(1-p).
\]
在 $0<p<1$ 上,
\[
 h'(p)=\log_2\frac{1-p}{p},\qquad
 h''(p)=-\frac1{\ln2}\left(\frac1p+\frac1{1-p}\right)<0.
\]
所以 $h$ 在 $p=1/2$ 取得最大值 $1$, 在两端为零.
若 $G$ 的质量函数为 $p(1-p)^{j-1}$, $j\geq1$, $0<p<1$, 则
\begin{align*}
 H(G)&=-\sum_{j\geq1}p(1-p)^{j-1}\bigl[\log_2p+(j-1)\log_2(1-p)\bigr]\\
 &=-\log_2p-\frac{1-p}{p}\log_2(1-p)=\frac{h(p)}p.
\end{align*}
这里使用 $\E(G-1)=(1-p)/p$. 特别地 $p=1/2$ 时 $H(G)=2$, 虽然可能取值有无穷多个, 平均信息量仍有限.
\end{solution}

\begin{proposition}[熵的凹性]
在有限概率单纯形 $\{(p_1,\ldots,p_m):p_i\geq0,\sum p_i=1\}$ 上,
$H(p)=-\sum_i p_i\log_2p_i$ 为严格凹函数. 即对两个不同概率向量 $p,q$ 及 $0<t<1$,
\[
 H(tp+(1-t)q)>tH(p)+(1-t)H(q).
\]
\end{proposition}
\begin{proof}
令 $g(s)=-s\log_2s$, $g(0)=0$. 在 $(0,1]$ 上 $g''(s)=-1/(s\ln2)<0$, 因而任意两个不同正数之间严格凹. 若一个端点为零, 另一个为 $a>0$, 则
\[
 g(ta)-tg(a)=-ta\log_2t>0\quad(0<t<1),
\]
仍有严格凹性. 对每个坐标应用凹性再相加得到不等式. $p\ne q$ 时至少一个坐标不同, 该坐标的不等式严格, 因而总和严格.
\end{proof}
凹性的意思是: 先随机选择一种分布, 又不告诉观察者选了哪一种, 混合分布的熵至少等于选择已知时的平均熵. 后面将以此证明条件化平均减少不确定性.

\originalsection{相对熵与最大熵}
比较两个概率分布需要一个非负量. 我们先证明它, 再用它给出最大熵和编码下界, 这样各结论有同一个明确的依据.
\begin{definition}[相对熵]
对同一有限集合上的分布 $p,q$, 定义
\[
 D(p\Vert q)=\sum_{i:p_i>0}p_i\log_2\frac{p_i}{q_i}.
\]
若某个 $p_i>0$ 而 $q_i=0$, 则定义 $D(p\Vert q)=+\infty$.
\end{definition}
相对熵一般不对称, 因而不是距离. 它衡量的是使用 $q$ 作为概率描述, 与实际分布 $p$ 不符所带来的平均对数损失.

\begin{lemma}[相对熵非负]
$D(p\Vert q)\geq0$, 等号成立当且仅当 $p=q$.
\end{lemma}
\begin{proof}
无穷的情况已经成立. 否则对所有 $p_i>0$ 使用 $\ln z\leq z-1$, 得到
\[
 (\ln2)D(p\Vert q)
 =-\sum_{p_i>0}p_i\ln\frac{q_i}{p_i}
 \geq\sum_{p_i>0}(p_i-q_i)
 =1-\sum_{p_i>0}q_i\geq0.
\]
不等式 $\ln z\leq z-1$ 可由 $z-1-\ln z$ 在 $z=1$ 取得最小值零证明; 等号只在 $z=1$ 成立. 所以总等号要求每个正概率坐标满足 $q_i=p_i$, 并且 $q$ 在其余坐标上总质量为零, 即 $p=q$. 反方向直接代入.
\end{proof}

\begin{theorem}[有限集合上的最大熵]
若 $X$ 至多有 $m$ 个可能取值, 则
\[
 0\leq H(X)\leq\log_2m.
\]
左等号当且仅当 $X$ 几乎必然为常量; 右等号当且仅当在指定的 $m$ 个取值上均匀分布.
\end{theorem}
\begin{proof}
各项 $-p_i\log_2p_i\geq0$. 若其和为零, 每个正 $p_i$ 必须为 $1$, 因而只有一个正概率取值. 取 $q_i=1/m$, 有
\[
 D(p\Vert q)=\sum_i p_i\log_2p_i+\log_2m=\log_2m-H(p).
\]
相对熵非负和等号条件分别给出右侧不等式及其等号条件.
\end{proof}
“有 $m$ 个可能结果”与“有 $\log_2m$ 比特的不确定性”只有在均匀分布时相同. 一个结果概率接近 $1$ 时, 即使其余结果很多, 熵也可能很小.

\originalsection{联合熵、条件熵与链式法则}
\begin{definition}
设 $X,Y$ 取有限多个值, 联合质量函数为 $p(x,y)$, 边缘质量函数为 $p_X,p_Y$. 定义
\[
 H(X,Y)=-\sum_{x,y}p(x,y)\log_2p(x,y).
\]
对 $p_Y(y)>0$, 定义
\[
 H(X\mid Y=y)=-\sum_xp(x\mid y)\log_2p(x\mid y),
 \qquad
 H(X\mid Y)=\sum_yp_Y(y)H(X\mid Y=y).
\]
零概率条件值在平均中不作贡献, 无需给它指定条件分布.
\end{definition}
$H(X\mid Y=y)$ 随观察值 $y$ 改变, 而 $H(X\mid Y)$ 是观察 $Y$ 前计算的平均剩余信息量, 是一个数, 不是条件期望随机变量.

\begin{theorem}[熵的链式法则]
\[
 H(X,Y)=H(Y)+H(X\mid Y)=H(X)+H(Y\mid X).
\]
更一般地,
\[
 H(X_1,\ldots,X_n)=H(X_1)+\sum_{j=2}^nH(X_j\mid X_1,\ldots,X_{j-1}).
\]
\end{theorem}
\begin{proof}
对 $p(x,y)>0$, 因 $p(x,y)=p_Y(y)p(x\mid y)$,
\begin{align*}
 H(X,Y)&=-\sum_{x,y}p(x,y)\log_2p_Y(y)
         -\sum_{x,y}p(x,y)\log_2p(x\mid y)\\
 &=-\sum_yp_Y(y)\log_2p_Y(y)
   +\sum_yp_Y(y)\left[-\sum_xp(x\mid y)\log_2p(x\mid y)\right].
\end{align*}
有限和允许改变求和次序; 零概率项按定义为零. 这就是第一式. 交换 $X,Y$ 得到第二式. 对向量 $(X_1,\ldots,X_{n-1})$ 和 $X_n$ 使用两变量公式并归纳, 得到多变量公式.
\end{proof}

\begin{theorem}[条件化减少平均熵]
\[
 0\leq H(X\mid Y)\leq H(X),\qquad H(X,Y)\leq H(X)+H(Y).
\]
$H(X\mid Y)=H(X)$ 当且仅当 $X,Y$ 独立. $H(X\mid Y)=0$ 当且仅当存在函数 $f$ 使 $X=f(Y)$ 几乎必然成立.
\end{theorem}
\begin{proof}
令 $v_y=(p(x_1\mid y),\ldots,p(x_m\mid y))$. 边缘概率向量为
$v=\sum_y p_Y(y)v_y$. 熵的凹性给出
\[
 H(X)=H\left(\sum_yp_Y(y)v_y\right)
 \geq\sum_yp_Y(y)H(v_y)=H(X\mid Y).
\]
严格凹性表明等号要求所有具有正权重的 $v_y$ 相同, 从而它们都等于 $v$. 这正是 $p(x\mid y)=p_X(x)$, 等价于独立. 也可直接计算
\[
 H(X)-H(X\mid Y)=D(p_{X,Y}\Vert p_Xp_Y)\geq0,
\]
其中任何 $p(x,y)>0$ 的点两边缘概率都正, 不会出现分母为零.
非负性来自定义; 联合熵上界来自链式法则. 平均条件熵为零时, 每个正概率 $y$ 上的条件熵都为零, 该条件分布集中在唯一的 $f(y)$ 上. 反之若 $X=f(Y)$, 这些条件分布都集中在单点, 条件熵为零.
\end{proof}
这个定理是平均意义的陈述. 某个特定观察值可能使熵增加. 例如 $Y$ 以 $9/10$ 的概率等于 $0$, 此时 $X=0$; $Y=1$ 时 $X$ 在 $1,2$ 上各占一半. 则 $H(X\mid Y=1)=1$, 而 $H(X)=h(1/10)+1/10<1$. 罕见观察揭示“目前进入了更不确定的情形”, 与平均条件熵 $H(X\mid Y)=1/10$ 较小并不矛盾.

\begin{definition}[互信息]
$ I(X;Y)=H(X)+H(Y)-H(X,Y)=H(X)-H(X\mid Y)$.
\end{definition}
链式法则说明它对称; 上一定理说明它非负, 且为零恰好是独立. 当 $Y=f(X)$ 时, $I(X;Y)=H(Y)$, 因为获知 $X$ 已经确定了 $Y$.

\begin{example}
$X$ 在 $\{0,1,2,3\}$ 上均匀分布, $Y$ 是 $X$ 的奇偶性. 计算联合熵、条件熵和互信息.
\end{example}
\begin{solution}
$H(X)=2$, $Y$ 为公平二元变量, $H(Y)=1$. 给定奇偶性后仍有两个等可能 $X$ 值, 所以 $H(X\mid Y)=1$. 给定 $X$ 后 $Y$ 确定, 故 $H(Y\mid X)=0$. 链式法则给出 $H(X,Y)=2$, 互信息为 $1$. 观察奇偶性正好消除一个比特的不确定性.
\end{solution}

\originalsection{二叉决策树与前缀码}
一个是非问题给出一个二进制位. 接下来问什么可以依赖已有回答, 所以完整策略是一棵二叉树: 根是开始, 内部结点是问题, 两条边是两个回答, 叶子是已经确定的结果. 某个结果所需问题数就是对应叶子的深度.

\begin{definition}[前缀码]
为每个可能结果 $x_i$ 指定有限二进制串 $c_i$. 若任一个码字都不是另一个码字的前缀, 则称为前缀码. 其码长为 $\ell_i$, 平均码长为 $L=\sum_ip_i\ell_i$.
\end{definition}
前缀条件允许即时解码: 从左到右走二叉树, 一到叶子就读出一个结果并回到根, 无需额外分隔符. 只有一个可能结果时允许空码字, 长度为零. 一般的唯一可解码码不必是前缀码; 本章证明针对与是非问题策略直接对应的前缀码.

\begin{lemma}[Kraft 不等式及其逆命题]
有限前缀码的长度满足 $\sum_i2^{-\ell_i}\leq1$. 反之, 非负整数 $\ell_i$ 若满足此不等式, 就存在具有这些长度的前缀码.
\end{lemma}
\begin{proof}
令 $L_*=\max_i\ell_i$, 考虑所有长 $L_*$ 的二进制串, 共 $2^{L_*}$ 个. 以 $c_i$ 为开头的串有 $2^{L_*-\ell_i}$ 个. 前缀条件保证这些集合不相交, 因而相加不超过 $2^{L_*}$; 除以它便得不等式.

逆命题将长度从小到大排列. 在深度 $\ell_i$ 处, 每个已选的深度 $\ell_j\leq\ell_i$ 的叶子禁用 $2^{\ell_i-\ell_j}$ 个位置. 因而选择第 $i$ 个码字时已有禁用位置数为
$\sum_{j<i}2^{\ell_i-\ell_j}$. 由 $\sum_{j\leq i}2^{-\ell_j}\leq1$ 得
\[
 \sum_{j<i}2^{\ell_i-\ell_j}\leq2^{\ell_i}-1.
\]
所以至少有一个位置可选. 选定它, 并在之后禁用其后代. 因选择按深度递增, 新叶子不会成为旧叶子的祖先, 也不是旧叶子的后代. 依次完成, 得到所需前缀码.
\end{proof}

\begin{theorem}[无噪声编码界]
对有限分布 $p$, 任意前缀码的平均长度满足 $L\geq H(p)$. 存在前缀码使
\[
 H(p)\leq L<H(p)+1.
\]
当所有正概率均为 $2$ 的负整数次幂时, 可以取得 $L=H(p)$.
\end{theorem}
\begin{proof}
只编码正概率结果. 对给定码令 $a=\sum_i2^{-\ell_i}\leq1$, $q_i=2^{-\ell_i}/a$. 相对熵非负给出
\[
 D(p\Vert q)=\sum_ip_i\log_2p_i+\sum_ip_i\ell_i+\log_2a
 =L-H(p)+\log_2a\geq0.
\]
所以 $L\geq H(p)-\log_2a\geq H(p)$.

为构造上界, 取 $\ell_i=\lceil-\log_2p_i\rceil$. 因为
$2^{-\ell_i}\leq p_i$, Kraft 逆命题保证这些长度可实现. 又因
$-\log_2p_i\leq\ell_i<-\log_2p_i+1$, 乘以 $p_i$ 后求和得所需界.
若 $p_i=2^{-k_i}$, 取 $\ell_i=k_i$, 则 Kraft 和为 $1$, 且 $L=\sum_ip_ik_i=H(p)$.
\end{proof}
下界证明还给出等号条件: $a=1$ 且 $p_i=q_i$, 即 $p_i=2^{-\ell_i}$. 在决策树中, 这意味着每个被访问的分叉都把余下概率分成两半. 一般概率没有整数比特数, 因而不能期待单个符号编码恰好达到熵.

\begin{example}
一个传感器输出四种状态, 概率为 $(1/2,1/4,1/8,1/8)$. 构造达到熵下界的编码, 并解释如何解码.
\end{example}
\begin{solution}
四种状态记为 $a,b,c,d$, 选码字 $0,10,110,111$. 码长分别为 $1,2,3,3$, 没有码字是另一个的前缀. 平均长度
\[
 L=\frac12+\frac24+\frac38+\frac38=\frac74=H(X).
\]
例如收到串 $101110110$, 从根逐位读得 $10|111|0|110$, 即 $b,d,a,c$. 竖线是解码后标出的边界, 传输时不需要发送它们.
\end{solution}

\originalsection{Huffman 编码与分组压缩}
上述向上取整法很容易实现, 但不总是最优. Huffman 算法反复合并最小的两个概率, 最后从合并树读出二进制码. 为什么可以安全地把最小概率放在最深处, 需要一个交换论证.

\begin{proposition}[Huffman 算法的最优性]
对有限个正概率结果, 反复将两个最小概率 $p_a,p_b$ 合并为一个概率 $p_a+p_b$, 对缩小后的集合编码, 再把合并结果的码字后添 $0,1$, 得到平均码长最小的前缀码.
\end{proposition}
\begin{proof}
只有一个结果时结论成立. 对至少两个结果, 最优树可取为每个内部结点都有两个孩子的树: 若只有一个孩子, 删除该无效分叉会缩短其下所有叶子的码长. 这种树最多有 $m-1$ 个内部结点, 所以深度不超过 $m-1$, 可能树及叶子标号只有有限多个, 最小值确实存在.

最深的一个叶子的兄弟也是叶子, 否则兄弟还有更深的后代. 若 $p_i>p_j$ 而 $\ell_i>\ell_j$, 交换这两个叶子的标签会使平均码长改变
\[
 p_i\ell_j+p_j\ell_i-p_i\ell_i-p_j\ell_j
 =-(p_i-p_j)(\ell_i-\ell_j)<0.
\]
因此可以让两个最小概率占据最深的一对兄弟叶子, 而不增加平均码长. 将这对兄弟合并, 新树平均长度比原树减少 $p_a+p_b$. 合并后的树必须最优, 否则用更好的小树替换再拆开会改善原树. 归纳假设给出缩小后 Huffman 树最优, 拆开时增加的 $p_a+p_b$ 与小树选择无关, 故所得原树最优.
\end{proof}

\begin{example}
对分布 $(0.4,0.3,0.2,0.1)$ 求一个最优前缀码.
\end{example}
\begin{solution}
先合并 $0.1,0.2$ 为 $0.3$, 再将两个 $0.3$ 合并为 $0.6$, 最后合并 $0.4,0.6$. 可以取码字 $0,10,110,111$, 分别对应 $0.4,0.3,0.2,0.1$. 平均长度为
$0.4+2(0.3)+3(0.2)+3(0.1)=1.9$.
熵为 $-\sum p_i\log_2p_i\approx1.8464$, 小于 $1.9$. 两个 $0.3$ 的处理顺序及每次左右分支可以互换, 最优码未必唯一.
\end{solution}

\begin{corollary}[分组编码接近熵率]
若 $X_1,\ldots,X_n$ 独立同分布, 每个变量取有限多个值, 则存在对整个向量的一次前缀编码, 平均总长度 $L_n$ 满足
\[
 nH(X_1)\leq L_n<nH(X_1)+1,
 \qquad \frac{L_n}{n}\longrightarrow H(X_1).
\]
\end{corollary}
\begin{proof}
独立性使每个条件熵等于无条件熵, 链式法则给出
$H(X_1,\ldots,X_n)=nH(X_1)$. 将整个向量当作一个具有有限个可能取值的随机变量应用编码定理即可. 最后的一个比特误差摊到 $n$ 个符号上成为 $1/n$.
\end{proof}
这解释了分组压缩的收益: 单个符号的码长必须是整数, 但一组符号可以共同承担取整误差. 这一定理是平均长度结果, 并不保证每一条输入序列都得到短码; 罕见序列仍可能很长.

\originalsection{练习与解答}
\begin{exercise}
$X$ 为公平比特, $B$ 为独立的伯努利 $(q)$ 比特, $Y=X+B\pmod2$. 求 $H(Y)$、$H(Y\mid X)$、$I(X;Y)$, 并解释 $q=0,1/2,1$.
\end{exercise}
\begin{solution}
$\Prob(Y=1)=\tfrac12q+\tfrac12(1-q)=1/2$, 故 $H(Y)=1$. 给定 $X$, $Y$ 的条件分布是 $(q,1-q)$ 或 $(1-q,q)$, 两者熵均为 $h(q)$. 因而 $H(Y\mid X)=h(q)$, $I(X;Y)=1-h(q)$. $q=0$ 时 $Y=X$, $q=1$ 时 $Y=1-X$, 二者均完全确定 $X$, 互信息为 $1$. $q=1/2$ 时输出不含关于输入的信息, 两者独立, 互信息为零.
\end{solution}

\begin{exercise}
三种输出的概率为 $(1/2,1/3,1/6)$. 求最优前缀码的平均长度, 并比较单符号编码与 $n$ 符号分组编码的每符号误差上界.
\end{exercise}
\begin{solution}
Huffman 算法合并 $1/3,1/6$, 得到两叶概率 $1/2,1/2$. 取码字 $0,10,11$, 平均长度为 $3/2$. 熵为
\[
 H=\frac12+\frac13\log_23+\frac16\log_26
 =\frac23+\frac12\log_23\approx1.4591.
\]
单符号最优码多花约 $0.0409$ 比特; 一般分组上界则保证 $L_n/n-H<1/n$. 分组编码要处理更多码字, 所以接近熵率同时带来存储和延迟方面的代价.
\end{solution}

\begin{exercise}
证明 $H(X\mid Y,Z)\leq H(X\mid Y)$, 并给出等号条件.
\end{exercise}
\begin{solution}
在每个正概率 $Y=y$ 的条件空间上应用条件化减少熵定理,
\[
 H(X\mid Y=y)\geq\sum_z\Prob(Z=z\mid Y=y)H(X\mid Y=y,Z=z).
\]
乘以 $\Prob(Y=y)$ 求和得到所需不等式. 等号成立当且仅当对每个正概率 $y$, 在该条件分布下 $X,Z$ 独立, 即
$p(x,z\mid y)=p(x\mid y)p(z\mid y)$.
\end{solution}
